\scanpage{75}
\mspage{- 19 -}

On notera que tout système de $q$ séries entières $u_i$ sans terme constant $(1\leq i\leq q)$ sommables dans un voisinage de l'origine, est substituable da ns chacune des séries entières $\handwritten{\sum_{(n_j)}}h_{i;n_1..n_q}y_1^{n_1}..y_q^{n_q}$ (cor. du th.2) ; si le problème a une solution, les séries entières (par rapport aux $x_i$) des deux membres de (5) \unsure{sort|sont} donc sommables au voisinage de l'origine, et par suite doivent a\unsure{-|{}}voir leurs coefficients \uline{identiques}.

Supposons donc le problème résolu, et désignons par $c_i(n_1,..,{}_{q};r_1,..,r_p)$ le \unsure{cœfficient|coefficient} de $x_1^{r_1}\ldots x_p^{r_p}$ dans la série $h_{i;n_1..n_q}$, par $c_i(r_1,..,r_p)$ le coefficient du même terme dans $h_{i0}$ par $d_i(m_1,..,m_p)$ le coefficient de $x_1^{m_1}..x_p^{m_p}$ dans $u_i$. L'identification des séries des deux membres de (5) donne les conditions

\[
\left\{\begin{aligned}
d_i(0,0,\ldots,0)&=0\\
d_i(m_1,..,m_p)&=c_i(m_1,..,m_p)+\\
&\quad+\handwritten{\sum_{(n_j)(r_i)}}(c_i(n_1,..,n_q;r_1,..,r_p)
\handwritten{\prod_{j=1}^q\prod_{k=1}^{n_j}}d_j(t_{kj1},t_{kj2},\ldots,t_{kjp}))
\end{aligned}\right.\tag*{(6)}
\]

la sommation du second membre étant étendue à tous les systèmes d'entiers $\geq 0$, $n_j,r_i$ et $t_{kji}$ satisfaisant aux conditions $\handwritten{\sum_{i=1}^{p}}t_{kji}\geq 1$, $\handwritten{\sum_{j=1}^q}n_j\geq 2$ et

\[
m_i=r_i+\handwritten{\sum_{j=1}^{q}\sum_{k=1}^{n_j}}t_{kji}\qquad(1\leq i\leq p)\ .\tag*{(7)}
\]

On tire de ces conditions que

\[
\handwritten{\sum_{i=1}^{p}}m_i=\handwritten{\sum_{i=1}^{p}}r_i+
\handwritten{\sum_{j=1}^{q}\sum_{k=1}^{n_j}\sum_{i=1}^{p}}t_{kji}
\]

Si on tient compte de ce que $\handwritten{\sum_{i=1}^p}t_{kji}\geq 1$ et $\handwritten{\sum_{j=1}^{q}}n_j\geq 2$, on voit que pour tout système $(t_{kji})_{1\leq i\leq p}$ qui figure au second membre de (7), on a $\handwritten{\sum_{i=1}^p}t_{kji}<\handwritten{\sum_{i=1}^{p}}m_i$, en d'autres termes, les relations (6) déterminent les coefficients des termes de degré \uline{total} $s$ dans $u_i$ en fonction de ceux des termes de degré total $<s$ dans les $u_j$ $(1\leq j\leq q)$ ; par récurrence sur $s$, on voit donc que les équations (6) déterminent les coefficients
% archive sheet number: 74

